\section{Runtime, Latency, and Governance Evidence}
\label{app:prof-runtime-evidence}

This appendix section collects the evidence surfaces that justify the
runtime contract in the main IEEE draft.  The purpose is not to flood the
reader with every locked artifact in the repo.  It is to show that the
adapter, benchmark, and CertOS claims are backed by measured runtime
latency, explicit replay tables, and a visible governance gate.

\subsection{Latency and Footprint}
The latency story is straightforward.  The runtime chapter
(\texttt{chapters/ch22\_latency\_systems\_footprint.tex}) measures the DC3S
pipeline with \texttt{time.perf\_counter\_ns()} and reports sub-millisecond
step time on the locked battery dispatch surface.  The key component
latencies are the reliability scorer, the drift updater, the uncertainty
set builder, and the action repair stage.

\begin{table}[t]
\centering
\caption{Per-component DC3S latency summary from the locked runtime study.}
\label{tab:prof-latency-summary}
\begin{tabular}{lrrrrr}
\toprule
\textbf{Component} & \textbf{Mean} & \textbf{Median} & \textbf{P95} & \textbf{P99} & \textbf{Max} \\
\midrule
Reliability scoring (OQE)   & 0.026 & 0.021 & 0.043 & 0.130 & 0.712 \\
Drift update (Page-Hinkley) & 0.000 & 0.000 & 0.001 & 0.001 & 0.160 \\
Uncertainty set build (RUI) & 0.012 & 0.011 & 0.011 & 0.047 & 0.724 \\
Action repair (SAF/Shield)  & 0.002 & 0.002 & 0.002 & 0.003 & 0.210 \\
\midrule
\textbf{Full DC3S step}     & \textbf{0.042} & \textbf{0.037} & \textbf{0.041} & \textbf{0.193} & \textbf{2.416} \\
\bottomrule
\end{tabular}
\end{table}

The same section also provides the best compact footprint statement for a
professor-facing draft: the full step fits comfortably inside a 30-second
dispatch cycle, and the runtime memory proxy remains dominated by the
forecast and residual buffers rather than by certificate bookkeeping.
That makes the runtime layer operationally cheap enough to be believable
as a Physical AI safety wrapper.

\subsection{Runtime Budget Across Domains}
The cross-domain runtime-budget matrix in
\texttt{reports/publication/tbl\_orius\_runtime\_budget\_matrix.tex} is the
best evidence that the runtime contract is not battery-specific.  It shows
that the active three-domain program preserves sub-millisecond P95 latency while
keeping certificate rates at 100 percent on all promoted rows.

\begin{table*}[t]
\centering
\scriptsize
\caption{Runtime and governance summary for the professor-facing ORIUS draft.}
\label{tab:prof-runtime-governance}
\begin{tabularx}{\textwidth}{@{}p{0.24\textwidth}p{0.16\textwidth}p{0.14\textwidth}p{0.14\textwidth}p{0.12\textwidth}p{0.14\textwidth}@{}}
\toprule
\textbf{Domain} & \textbf{Row} & \textbf{P95 latency (ms)} & \textbf{Fallback (\%)} & \textbf{Audit / lifecycle (\%)} & \textbf{Boundary} \\
\midrule
Battery Energy Storage & Witness row & 0.812 & 100.0 & 100.0 & deepest theorem-to-artifact closure\\
Autonomous Vehicles & Defended bounded row & 0.744 & 100.0 & 100.0 & bounded TTC plus predictive-entry-barrier contract\\
Medical and Healthcare Monitoring & Defended bounded row & 0.701 & 100.0 & 100.0 & bounded monitoring and alert-release semantics\\
\bottomrule
\end{tabularx}
\end{table*}


This table should stay close to the runtime section in the professor draft
because it directly supports the claim that the universal runtime layer is
lightweight across domains while still respecting the evidence tier.

\subsection{Governance Evidence}
The governance story is best captured by the locked lifecycle matrix in
\texttt{reports/publication/tbl\_orius\_governance\_lifecycle\_matrix.tex} and
the domain-closure matrix in \texttt{reports/publication/orius\_domain\_closure\_matrix.csv}.
Together they show that CertOS is not a decorative abstraction: the same
issue/validate/expire/fallback logic is observed in the defended rows.

\begin{table}[t]
\centering
\caption{CertOS runtime summary from the battery lifecycle demo.}
\label{tab:prof-certos-summary}
\begin{tabular}{lr}
\toprule
\textbf{Metric} & \textbf{Value} \\
\midrule
Total steps & 96 \\
Valid steps & 47 \\
Degraded steps & 16 \\
Fallback steps & 33 \\
Runtime interventions & 33 \\
Audit entries & 96 \\
Audit completeness & 100.0\% \\
Step latency & 0.59 ms \\
\bottomrule
\end{tabular}
\end{table}

The appendix should end with the governance boundary spelled out plainly:
the runtime layer is fully visible, the certificates are auditable, and the
fallback path is explicit, but the claim remains tiered by domain. Battery
is the reference witness; AV and healthcare are defended bounded rows. That
is enough for a strong IEEE draft, and it is more defensible than flattening
the evidence into one universal-number story.

\subsection{Repository maturity visual}
Appendix~B is also the right place for one explicit repository-level visual.
The professor draft should show the reader that ORIUS is not only a narrative
proposal but a maintained research system with a real public code surface and
a generated architecture map.

\begin{figure}[t]
\centering
\includegraphics[width=\linewidth]{orius_repo_github_overview}
\caption{GitHub-style repository overview used as a maturity visual in the
professor draft. This figure is supplemental and does not carry scientific
claims by itself.}
\label{fig:prof-github-overview}
\end{figure}

\begin{figure}[t]
\centering
\includegraphics[width=\linewidth]{fig01_architecture}
\caption{Repository-backed architecture visual reused in Appendix~B to connect
the code surface to the runtime stack argued in the main paper.}
\label{fig:prof-appendix-architecture}
\end{figure}
